\section{History events and conditional estimates}
\label{frames:chapter}

The preceding sections obtained marginal estimates both for every fixed
input list and after averaging the list. We now retain information about
trajectory histories before passing to endpoint laws. For a fixed
partition and fixed orders on its complements, one event gives pointwise
coset caps after its energy constraints are removed in descending order.
A stopped estimate for a fixed list charges imbalance only once across
the list. A third construction gives an entropy comparison for every
probability law supported on one common event of cyclic histories.

The frame comparison, the balance estimate, and the retained conclusion
play separate roles in these arguments. In particular, the deferred
shear frame used in the final application preserves the distance for
each fixed input list; that application averages the labels to obtain
simultaneous balance across affine half-spaces. We give the three history
constructions and their full-deck consequences first, then return to
that averaged application.

Throughout this section $n=2^d$, $V=\mathbb F_2^d$, and
$G=\operatorname{Sym}(V)$. Permutations send input positions to output
positions and multiply by composition. Write $U_G$ for uniform measure.
For $v\ne0$, retain the matching subgroup $C_v$: it consists of the
permutations preserving every pair $\{x,x+v\}$ setwise. A uniform element
of $C_v$ is a layer of independent fair switches. Removing the
accumulated coordinate rotation from the physical
shuffle gives the cyclic directions $w_t=e_{1+((t-1)\bmod d)}$.
In particular a block of length divisible by $d$ has exactly the same
endpoint permutation in these coordinates as in physical coordinates.

\begin{remark}\label{frames:lower-import}
Every upper bound in this section is paired with the support estimate of
Lemma~\ref{lem:lower}: there are at most $2^{tn/2}$ outcomes after $t$
physical shuffles, so distance from uniform is at least
$1-2^{tn/2}/n!$. For integer $t\le d$ this tends uniformly to one.
Consequently each upper bound below has order $d$, and the mixing time
at threshold $1/4$ is at least $d$ for sufficiently large $d$.
\end{remark}

\subsection{Preserving the matching of labels}
\label{frames:invariance-section}

\begin{lemma}[Matching and input-list invariance]
\label{frames:matching-invariance}
Fix an integer $T\ge d$ and directions $v_1,\ldots,v_T$ such that
every consecutive $d$ form a
basis. There are invertible linear maps $L_0,\ldots,L_T$ with
$L_0e_i=v_i$ for $1\le i\le d$ such that, if $g_t$ is the cyclic-layer
process, then $L_tg_tL_0^{-1}$ is the independent switch process with
directions $v_t$. This assertion is an equality in law conditional on
the entire fixed sequence. If $L_0=I$, the matching on trajectory labels
before each switch is unchanged. For general $L_0$, the same statement
holds after relabeling the initial labels by $L_0$. Consequently the
average, over uniform ordered input lists of any fixed length, of the
endpoint distance to a uniform injection is unchanged.
\end{lemma}

\begin{proof}
Lemma~\ref{shear:lem:general-frame} gives exactly the asserted path-law
identity, with $L_0e_i=v_i$. Its construction applies to every fixed
schedule with consecutive basis windows and permits an arbitrary
initial basis. It remains to check what happens to trajectory labels.

If labels $a,b$ are paired in the cyclic process just before layer $t$,
then $g_{t-1}(b)=g_{t-1}(a)+w_t$. The column identity in that frame
construction is $L_{t-1}w_t=v_t$. Hence their transformed positions,
with initial labels changed by $L_0$, differ by $v_t$ and are paired
in the transformed process. The implication reverses because
$L_{t-1}$ is invertible. When $L_0=I$, the labels themselves do not
change. For general $L_0$, its bijection preserves the uniform law
of the input list. The terminal bijection $L_T$ preserves uniform
injections and total variation. This proves both assertions.
\end{proof}

The maps $L_t$ may depend on future directions. We use them only after
fixing the complete schedule. In the probabilistic calculations below,
the conditional switch law given that schedule is the product of the
uniform matching laws. We reveal directions chronologically and average
them only in the prefix fields specified below.

\begin{lemma}[Conditional mass on unobserved positions]
\label{frames:holes}
Fix an integer $T\ge1$ and nonzero directions $v_1,\ldots,v_T$. Let $Z_0=I$ and
$Z_t=\xi_tZ_{t-1}$, where the $\xi_t$ are independent uniform switches
in directions $v_t$. Fix complete feasible trajectories of some observed
labels. If $m\ge1$ labels remain unobserved, let $F_t$ be the positions
not occupied by the observed labels at time $t$. For a specified
unobserved label, let $p_t$ be its conditional position law given those
complete observed trajectories.

The vector $p_t$ evolves from its initial point mass by averaging on
pairs with two positions in $F_{t-1}$ and deterministic transport on
pairs with one such position. It depends only on the directions and
observed paths through time $t$. The uniform vector $1/m$ is transported
to the uniform vector on $F_t$. Thus, with
$E_t=\sum_{x\in F_t}(p_t(x)-1/m)^2$, one has $0\le E_t\le1$ and
\begin{equation}
 E_{t-1}-E_t=\frac12\sum_{\{x,y\}\subset F_{t-1}\text{ paired at }t}
                    (p_{t-1}(x)-p_{t-1}(y))^2.
 \label{frames:hole-loss}
\end{equation}
\end{lemma}

\begin{proof}
The online averaging and transport rules are those of
Lemma~\ref{shear:lem:energy}; here we must also justify conditioning
on complete observed paths. A feasible specification of those paths
prescribes precisely the coins on the visited time-edge pairs.
Necessity is immediate. Conversely, those prescribed coins force the
paths, by induction over the layers. The path event is therefore a
cylinder in the independent coin array. All other coins remain
independent and fair, including unvisited coins at earlier times.

On an edge used by an observed label the coin is fixed; on an edge
with two unobserved positions it remains fair. Future path constraints
use later time-edge pairs and do not alter this earlier recursion.
Thus the conditional vector equals the same forward averaging and
transport flow obtained from the observed prefix. Applying the
exact decrement in Lemma~\ref{shear:lem:energy}, with its available
set equal to $F_{t-1}$, proves the displayed identity and transports
the uniform comparison vector as stated.
\end{proof}

The next consequence records the two conditioning fields used below.
The complete field is convenient for endpoint comparisons; the smaller
field is the one in which the next direction is averaged.

\begin{corollary}[Complete and chronological conditioning]
\label{frames:conditional-interface}
Fix integers $T\ge1$ and $1\le q\le n$. Let
$\mathbf v=(v_1,\ldots,v_T)$ be a random schedule of nonzero directions.
Set $Z_0=I$, and suppose that, conditional on $\mathbf v$, the
increments of $Z$ are independent uniform switches in directions
$v_1,\ldots,v_T$. Let $C=(c_1,\ldots,c_q)$ be a deterministic list of
distinct labels, or a random such list independent of the joint process
$(\mathbf v,(Z_t)_{t=0}^T)$. For $1\le j\le q$, put
\begin{align}
 \mathcal F_t^{C,j}
 &=\sigma\bigl(C,v_1,\ldots,v_t,\,
       Z_s(c_i):0\le s\le t,\ i<j\bigr), \label{frames:online-field}\\
 \mathcal H^{C,j}
 &=\sigma\bigl(C,\mathbf v,\,
       Z_s(c_i):0\le s\le T,\ i<j\bigr). \label{frames:terminal-field}
\end{align}
For $0\le t\le T$, define
$p_t(x)=\Pr\{Z_t(c_j)=x\mid\mathcal F_t^{C,j}\}$.
Then, for every $x\in V$,
\begin{equation}
 \Pr\{Z_t(c_j)=x\mid\mathcal H^{C,j}\}=p_t(x)
 \quad\text{almost surely}. \label{frames:full-path-prefix}
\end{equation}
For $t<T$, let
$K_t(v)=\Pr\{v_{t+1}=v\mid v_1,\ldots,v_t\}$. Then
\begin{equation}
 \Pr\{Z_t(c_j)=x,\ v_{t+1}=v\mid\mathcal F_t^{C,j}\}
 =p_t(x)K_t(v). \label{frames:direction-factorization}
\end{equation}
\end{corollary}

\begin{proof}
First fix the list, the complete schedule, and feasible complete
trajectories of the preceding labels. Lemma~\ref{frames:holes} identifies
the conditional mass on the left of \eqref{frames:full-path-prefix} with
the forward mass computed from the list, the direction prefix, and the
observed path prefix through $t$. It is therefore already
$\mathcal F_t^{C,j}$-measurable. The tower property, using
$\mathcal F_t^{C,j}\subseteq\mathcal H^{C,j}$, proves
\eqref{frames:full-path-prefix}.

Conditional on the complete schedule, the observed path prefix depends
only on the direction prefix, by the product-switch hypothesis. Since
the list is independent, its observation and that of the preceding
paths do not change the next-direction law $K_t$. Also, the field
generated by $\mathcal F_t^{C,j}$ and $v_{t+1}$ is contained in
$\mathcal H^{C,j}$. Taking the tower of
\eqref{frames:full-path-prefix} to this intermediate field shows that
the hidden mass remains $p_t$ after $v_{t+1}$ is revealed. Combining
these two statements proves \eqref{frames:direction-factorization}.
\end{proof}

In the shear construction the auxiliary data $\Lambda$ determine
$\mathbf v$ but may contain additional entries. Suppose, as in
Lemma~\ref{shear:lem:keys}, that the conditional law of the full virtual
trajectory given $\Lambda$ depends on $\Lambda$ only through
$\mathbf v$, and that a sampled list is independent of
$(\Lambda,(Z_t)_{t=0}^T)$. Conditional independence, followed by
conditioning on the observed paths, gives
\begin{equation}
 \Pr\{Z_t(c_j)=x\mid
       \mathcal H^{C,j}\vee\sigma(\Lambda)\}=p_t(x).
 \label{frames:terminal-disclosure}
\end{equation}
The averaging in \eqref{frames:direction-factorization} is performed
in $\mathcal F_t^{C,j}$ before this terminal disclosure. Once
$\Lambda$ is revealed at time $T$, the direction averaging is finished.
The separate conditional frame identity then identifies the mapped
endpoint law with the original coordinate chain.

\begin{lemma}[Fresh-direction contraction]
\label{frames:cross-energy}
Suppose, given the generated past, the next direction is uniform outside
a hyperplane $J\subset V$. Split the $m$ unobserved positions into the
two cosets, denoting the resulting sets by $F_0,F_1$ and their sizes
by $m_0,m_1$. For any centered real mass vector $z$
known from that past, the expected squared-norm loss obtained by
averaging on the paired unobserved positions is
\[
 \frac1n\sum_{x\in F_0,\ y\in F_1}(z_x-z_y)^2
 \ge \frac{\min(m_0,m_1)}n\sum_xz_x^2.
\]
If the next direction is uniform outside a smaller subspace contained
in $J$, the same lower bound holds with $1/(2n)$ in place of $1/n$.
\end{lemma}

\begin{proof}
Apply the energy calculation in Lemma~\ref{shear:lem:energy} to the
available set $F_0\cup F_1$ and the centered vector $z$. That proof
uses only $\sum_xz_x=0$, not positivity of a probability vector.
In its notation, take $D^0=F_0$ and $D^1=F_1$; the exact
cross-pair loss is the displayed sum divided by $n$.
When the forbidden subspace is smaller than $J$, every cross-coset
difference remains allowed and has probability at least $1/n$,
instead of $2/n$. Its averaging loss is still half the squared
difference, giving the denominator $2n$.
\end{proof}

\subsection{The Fourier input for the applications}
\label{frames:degree-section}

The remainder uses the general results of the companion
\emph{From partial permutation information to Fourier bounds}~\cite{partial-fourier}.
We record the normalizations that enter the numerical applications.
For every fixed $p>0$, Theorem~\Lref{l:degree-all-powers} gives
\[
 C_p=\sup_{s\ge1}\sum_{\lambda\vdash s}(f^\lambda)^{-p}<\infty,
 \qquad
 \sum_{\substack{\lambda\vdash s\\
                 \lambda\notin\{(s),(1^s)\}}}(f^\lambda)^{-p}\longrightarrow0.
\]
Here $f^\lambda$ is the degree of the irreducible representation of
$S_s$ indexed by $\lambda$. For a subprobability $\nu$ on $S_n$, use
$\widehat\nu(\rho)=\sum_g\nu(g)\rho(g)$, without dividing by its mass.
Let $r\ge1$, let $M_1,\ldots,M_r$ be disjoint nonempty subsets
of the labels, and let $K_i=\operatorname{Sym}(M_i)$ fix the complement
pointwise; the subsets need not cover all labels. If
$\nu(gK_i)\le K/[S_n:K_i]$, the disjoint-support orbit argument
in \eqref{shear:eq:noncovering-orbit-general}, which extends
\cite[Proposition~\Lref{l:orbit-span}]{partial-fourier}, gives
\begin{equation}
 \|\widehat\nu(\rho)\|_{\rm op}
 \le\sqrt{rKC_p}\,D^{-1/2+(p+2)/(2r)},\qquad D=\dim\rho.
 \label{frames:coset-fourier}
\end{equation}
Its one-vector orbit estimate is independent of representation
multiplicities. It applies to unequal blocks as well, although our
applications use equal ones.

We use the amplification consequence in Proposition~\Lref{l:amplification}:
if $0\le\nu\le\mu$, $\nu(S_n)\to1$, $\mu$ has zero sign mean, and
$\|\widehat\nu(\rho)\|_{\rm op}\le B D^{-\beta}$ for $D>1$,
where $B,\beta>0$ are fixed, then every fixed $k$ with
$2-2k\beta<0$ gives
$\|\mu^{*k}-U_G\|_{\rm TV}\to0$. It also applies when $\nu$ is a
probability at distance $o(1)$ from $\mu$ and has sign mean $o(1)$.
All norms use unnormalized Hilbert--Schmidt trace. We verify the sign,
mass, and physical block length at each application.

\subsection{A retained event for one partition}
\label{frames:pointwise-section}

Proposition~\ref{shear:prop:fixed-marginal} already gives a
total-variation bound for every fixed large input list. Here we first
choose a block partition and an order on each complement. We retain one
event on which the conditional energies for these particular orders
are simultaneously small. Removing their constraints in descending
order will give pointwise caps for every output of each complement:
a factor of two for the retained subprobability, or four after
normalization.

Fix a horizon $T\ge d$. The initial directions are fixed to $e_1,\ldots,e_d$, and thereafter
$v_t$ is uniform outside $J_t=\operatorname{span}(v_{t-d+1},\ldots,v_{t-1})$.
Write $\mathbf v=(v_1,\ldots,v_T)$, and let $Z_0=I$ and
$Z_t=\xi_t Z_{t-1}$ be the cumulative
variable-direction permutation, where, conditional on the directions,
the $\xi_t$ are independent uniforms in $C_{v_t}$. With the frame
starting at $L_0=I$, the cyclic endpoint has law $L_T^{-1}Z_T$.
The preceding transverse-layer estimate supplies the energy bound
needed for this simultaneous event.

\begin{lemma}[Balance supplied by an earlier crossing]
\label{frames:crossing-balance}
Fix integers $T\ge d$ and $r\ge1$. For any fixed observed labels and
any $d<t\le T$, if $m$ positions remain unobserved, then
\[
 \Pr\{\min(|F_{t-1}\cap J_t|,|F_{t-1}\setminus J_t|)<m/4\}
 \le2e^{-m/8}.
\]
For the hidden-card energy of Lemma~\ref{frames:holes}, interpreted in
this random process by Corollary~\ref{frames:conditional-interface},
if $m\ge n/r$ and $\alpha=1/(4r)$, then
\begin{equation}
 \mathbb E E_T\le(1-\alpha)^{T-d}+2e^{-n/(8r)}.
 \label{frames:crossing-recursion}
\end{equation}
\end{lemma}

\begin{proof}
The transverse-layer proof of Lemma~\ref{shear:lem:fixed-balance}
is uniform over the entire fixed direction schedule. It therefore
applies here even though the first basis is prescribed. Explicitly,
fix all directions and the observed-label configuration just before
layer $t-d$.
The direction at that layer crosses the cosets of $J_t$, because
$v_{t-d},\ldots,v_{t-1}$ form a basis. The following $d-1$ directions
lie in $J_t$ and preserve the two coset counts. If $a$ crossing pairs
contain two unobserved positions, one count after the crossing is
$a+\operatorname{Binomial}(m-2a,1/2)$. The conditional fair-coin law
is unchanged by fixing the directions. The binomial tail from that
lemma is at most $2e^{-m/8}$; if $m=2a$ the count is deterministic.
Removing this conditioning proves the stated unconditional bound.

Equation~\eqref{frames:full-path-prefix} makes the energy vector
measurable in the chronological field, and
\eqref{frames:direction-factorization} permits averaging the next
direction there. Lemma~\ref{frames:cross-energy} then contracts
energy by at least $\alpha=1/(4r)$ whenever the two counts are at
least $m/4$. We use the balance estimate only after taking
expectations. Since $0\le E\le1$,
\[
 \mathbb E E_t\le(1-\alpha)\mathbb E E_{t-1}
                    +2\alpha e^{-n/(8r)}.
\]
Iterate from $E_d\le1$ and sum the geometric series. This proves
\eqref{frames:crossing-recursion} without asserting concentration
conditional on the subsequently observed paths.
\end{proof}

\begin{lemma}[Removing exceptional histories in descending order]
\label{frames:backward-removal}
Let $T\ge d$ be an integer, and let $r\ge2$ divide $n$. Partition $V$
into equal input blocks $M_1,\ldots,M_r$ and, for each $i$, fix an order
$c_{i,1},\ldots,c_{i,q}$ on its complement, where $q=n-n/r$.
Write $\mathcal P$ for the partition together with these orders, and
let $K_i=\operatorname{Sym}(M_i)$ fix the complement pointwise.
Use the variable-direction process just defined, whose first frame is
$L_0=I$. For $1\le j\le q$, let
\[
 \mathcal H_{i,j}
 =\sigma\bigl(\mathbf v,\,
       Z_s(c_{i,h}):0\le s\le T,\ h<j\bigr),\qquad
 F_{i,j}=V\setminus\{Z_T(c_{i,h}):h<j\},
\]
and define
\[
 p_{i,j}(x)=\Pr\{Z_T(c_{i,j})=x\mid\mathcal H_{i,j}\},\qquad
 E_{i,j}=\sum_{x\in F_{i,j}}
       \left(p_{i,j}(x)-\frac1{n-j+1}\right)^2.
\]
Suppose $\mathbb E E_{i,j}\le\varepsilon_n$ uniformly. Put
$\mathcal G_{\mathcal P}=\bigcap_{i,j}\{E_{i,j}\le n^{-6}\}$, and let
$\mu$ be the original cyclic endpoint law at time $T$. Then
$\Pr(\mathcal G_{\mathcal P}^c)\le rq n^6\varepsilon_n$. The measure
\[
 \nu(g)=\Pr\{\mathcal G_{\mathcal P},\ L_T^{-1}Z_T=g\}
\]
satisfies $0\le\nu\le\mu$ and, for large $n$,
$\nu(gK_i)\le2U_G(gK_i)$ for every $g\in G$ and $1\le i\le r$. If
$a=\Pr(\mathcal G_{\mathcal P})=1-o(1)$, its normalized law
$\bar\nu=\nu/a$ satisfies $\bar\nu(gK_i)\le4U_G(gK_i)$ and
$\|\bar\nu-\mu\|_{\rm TV}\le1-a=o(1)$.
\end{lemma}

\begin{proof}
Markov's inequality and a union bound give the lost mass.
For each complete direction schedule the law of $L_T^{-1}Z_T$ is
$\mu$, so restricting to $\mathcal G_{\mathcal P}$ gives
$0\le\nu\le\mu$.
For an order, its $j$-th energy event is measurable from all directions
and the preceding paths. On it, every endpoint has conditional
probability at most
\[
 \frac1{n-j+1}+n^{-3}\le\frac{1+n^{-2}}{n-j+1}.
\]
To bound a specified list of distinct outputs intersected with
$\mathcal G_{\mathcal P}$,
first discard requirements belonging to other orders. Condition on all
directions and the first $q-1$ paths, bound the final target on its own
energy event, and then discard that event. Continue downwards through
the list. Each remaining target event and energy constraint is measurable
in the conditioning data at its turn. The target in variable coordinates
is known because $L_T$ is determined by all directions. We obtain
\[
 \Pr(\mathcal G_{\mathcal P},\text{specified outputs})
 \le\prod_{j=1}^q\frac{1+n^{-2}}{n-j+1}
 \le2\frac{(n-q)!}{n!}.
\]
These output specifications are precisely the left cosets $gK_i$.
Normalization costs $1/a$, at most two eventually. The total-variation
bound for $\bar\nu$ follows from Lemma~\ref{lem:subprobability-fourier}.
\end{proof}

\begin{proposition}[Two applications of the retained event]
\label{frames:pointwise-applications}
The physical shuffle has total-variation distance $o(1)$ from uniform
by time $4096d$. The choices $r=32$, $T=2561d$ give the separate
subprobability construction of Lemma~\ref{frames:backward-removal}, and
sixteen such blocks give distance $o(1)$ at time $40976d$.
\end{proposition}

\begin{proof}
Take $d\ge5$, which suffices for all block sizes in this proposition.
For the first assertion take $r=8$ and $T=512d$.
Equation~\eqref{frames:crossing-recursion} is at most $n^{-12}$ for
large $d$, since $(31/32)^{511d}\le e^{-511d/32}$.
The exceptional mass is at most $8n\,n^6n^{-12}=o(1)$.
Use the normalized law $\bar\nu$ in Lemma~\ref{frames:backward-removal} and the
direct inverse-square estimate of Lemma~\Lref{l:degree-inverse-square}.
Proposition~\Lref{l:orbit-span}, with $p=2$, gives
$\|\widehat{\bar\nu}(\rho)\|_{\rm op}\le\sqrt{32C_2}D^{-1/4}$.
Eight convolutions have Plancherel exponent $2-16/4=-2$.
The original block has unbiased sign because one independent switch
already has fair parity; the normalized modification has sign mean
$o(1)$. Proposition~\Lref{l:amplification} applies. Since $T=512d$
is a whole number of coordinate sweeps, eight physical blocks give
$8T=4096d$.

For the second construction take $r=32$ and $T=(1+80r)d$.
The first term in \eqref{frames:crossing-recursion} is at most
$e^{-20d}\le n^{-20}$, so the exceptional mass is at most
$32n^7(n^{-20}+2e^{-n/256})=o(1)$.
Keep the unnormalized subprobability, whose coset factor is $K=2$.
The paragraph ``Central-binomial powers eight and twenty'' following
Lemma~\Lref{l:degree-inverse-ten} in
Subsection~\Lref{l:degree-square-root-section} of~\cite{partial-fourier} gives
\[
 C_8=\sup_{s\ge1}\sum_{\lambda\vdash s}(f^\lambda)^{-8}<\infty,
 \qquad
 \sum_{\substack{\lambda\vdash s\\
                 \lambda\notin\{(s),(1^s)\}}}(f^\lambda)^{-8}
 \longrightarrow0\qquad(s\longrightarrow\infty).
\]
Applying Proposition~\Lref{l:orbit-span} of the same companion
with $p=8$, $K=2$, and $r=32$ gives exponent
$-1/2+10/64=-11/32$ and prefactor $\sqrt{64C_8}$.
Sixteen convolutions leave exponent $2-32(11/32)=-9$. The retained
measure is dominated by the physical block law, whose sign mean is zero.
Proposition~\Lref{l:amplification}, including its trivial and sign terms,
therefore applies. Here $T=2561d$ is again a whole number of sweeps, so
$16T=40976d$ physical shuffles suffice.
\end{proof}

The ordinary averaged-list route is given in
Proposition~\ref{frames:average-marginals}. It uses a random initial
basis and sampled labels, and retains the $b=64$, $T=2048d$ marginal
estimate and its $32768d$ full-deck application. The next construction
returns to a fixed input list and keeps a nested event through the
sequential comparison.

\subsection{A stopped estimate that pays for imbalance once}
\label{frames:stopped-section}

The earlier crossing-layer estimate gave exponential concentration.
The conditional-tuple viewpoint is related to the analysis of the
swap-or-not shuffle in \cite[Section~3]{hmr2014}; the schedule comparison
and estimates here are proved for the present shuffle.
There is another way to balance every fixed list: randomize the later
hyperplane while holding an earlier configuration fixed. Although this
only gives an inverse-linear tail, nesting the balance events lets us
pay this error once for the whole list.

\begin{proposition}[Stopped partial-observation estimate]
\label{frames:stopped-marginals}
Let $d\ge4$. Put $b=16$, $m=n/b$, $q=n-m$, and
$T=(2+20b)d=322d$.
For every fixed ordered list of $q$ distinct input positions, its image
under the physical shuffle at time $T$ has distance at most
\begin{equation}
 \delta_d=\frac{4T}{m}+\frac12n^{3/2}
                (1-1/(4b))^{(T-2d+1)/2}=o(1)
 \label{frames:stopped-delta}
\end{equation}
from a uniform injection, uniformly over the chosen list.
\end{proposition}

\begin{proof}
Start with the standard basis and generate directions by independent
fair coefficients
\[
 v_t=v_{t-d}+\sum_{t-d<s<t}c_{s,t}v_s\quad(t>d).
\]
Write $\mathbf c=(c_{s,t}:d<t\le T,\ t-d<s<t)$ for the complete
coefficient array. Conditional on this array, let
$Z_0=I$ and $Z_t=\zeta_tZ_{t-1}$, where the $\zeta_t$ are independent
uniform switches in directions $v_t$. This is the same conditional
uniform-complement schedule as before. For each fixed $\mathbf c$,
the frame comparison starts at $L_0=I$. Thus the distance for the
fixed endpoint list is unchanged by the output relabeling, and equals
the average over $\mathbf c$ of the conditional auxiliary distance.

For later use, let
\begin{equation}
 \mathcal A_h=\sigma\bigl(
       c_{s,u}:d<u\le h,\ u-d<s<u;\quad
       \zeta_u:1\le u\le h\bigr). \label{frames:stopped-generated-field}
\end{equation}
This field fixes the directions and the complete switch history through
layer $h$. The conditional law of these switches uses only the directions
through $h$, hence only coefficients with target time at most $h$.
Coefficients with later target times therefore remain independent fair
bits given $\mathcal A_h$.

Fix the input list $a_1,\ldots,a_q$, write $Y_t(j)=Z_t(a_j)$, and let
$F_{j,t}=V\setminus\{Y_t(i):i<j\}$, of size $n_j=n-j+1\ge m$.
Define the chronological field
\[
 \mathcal F_{j,t}
 =\sigma\bigl(v_1,\ldots,v_t,\,
       Y_s(i):0\le s\le t,\ i<j\bigr).
\]
The conditional path law given $\mathbf c$ depends only on its
direction sequence. Equations~\eqref{frames:full-path-prefix} and
\eqref{frames:terminal-disclosure}, with $\Lambda=\mathbf c$, therefore
identify the two versions of the hidden mass:
\[
 p_{j,t}(x)
 :=\Pr\{Y_t(j)=x\mid\mathbf c,\,
               Y_s(i):0\le s\le T,\ i<j\}
 =\Pr\{Y_t(j)=x\mid\mathcal F_{j,t}\}.
\]
Put
\[
 E_{j,t}=\sum_{x\in F_{j,t}}
             \left(p_{j,t}(x)-\frac1{n_j}\right)^2.
\]
Thus $E_{j,t}$ is measurable in $\mathcal F_{j,t}$, even though its
first definition uses the complete coefficient array and complete
preceding paths. For $t>d$ set
$J_t=\operatorname{span}(v_{t-d+1},\ldots,v_{t-1})$.
Say that step $(j,t)$ is balanced if both its cosets contain at least
$m/4$ members of $F_{j,t-1}$. Let $\mathcal B_j(t)$ require this for
all steps $2d,\ldots,t$, with no requirement when $t=2d-1$.
The event $\mathcal B_j(t)$ is $\mathcal F_{j,t-1}$-measurable, before
the direction at $t$ is revealed. Equation
\eqref{frames:direction-factorization} and
Lemma~\ref{frames:cross-energy}, followed by
$\mathcal B_j(t)\subseteq\mathcal B_j(t-1)$, give
\[
 \mathbb E[\mathbf1_{\mathcal B_j(t)}E_{j,t}]
 \le(1-1/(4b))\,
        \mathbb E[\mathbf1_{\mathcal B_j(t-1)}E_{j,t-1}].
\]
Thus, for $\mathcal B_j=\mathcal B_j(T)$,
\begin{equation}
 \mathbb E[\mathbf1_{\mathcal B_j}E_{j,T}]
 \le(1-1/(4b))^{T-2d+1}.
 \label{frames:stopped-energy}
\end{equation}
The events decrease as $j$ increases because the hole sets shrink and
the common threshold $m/4$ stays fixed.

The energy on the nested balance events is now controlled. It remains
to estimate the probability that the smallest reservoir fails balance;
this will be the only exceptional-event cost in the tuple comparison.
Fix $t\ge2d$ and write $h=t-d$.
Let $\ell_t$ be the unique linear form vanishing on $J_t$ and taking
value one on $v_h$. Given $\mathcal A_h$, this form
is uniform among the forms with $\ell_t(v_h)=1$.
Indeed use the past basis $v_{h-d+1},\ldots,v_h$. The equations
$\ell_t(v_{h+i})=0$, $1\le i<d$, give
\[
 \ell_t(v_{h-d+i})
 =\sum_{s=h-d+i+1}^{h}c_{s,h+i}\ell_t(v_s).
\]
Solve backwards from $i=d-1$ to $i=1$. At each stage the term
$c_{h,h+i}\ell_t(v_h)=c_{h,h+i}$ is a fresh fair bit, independent of
the coefficients at later stages, of the other coefficients in the
current equation, and of $\mathcal A_h$. This proves the
uniformity claim.

For $F=F_{q,h}$ put $Z=\sum_{x\in F}(-1)^{\ell_t(x)}$.
The set $F$ is $\mathcal A_h$-measurable, and
$\mathbb E[Z^2\mid\mathcal A_h]\le|F|=n_q$: diagonal terms are one, off-diagonal terms
with difference $v_h$ are minus one, and every other off-diagonal term
has mean zero. The directions at times $h+1,\ldots,t-1$ lie in $J_t$,
so the coset counts from the post-step-$h$ configuration to time $t-1$
are preserved. A count below $m/4$ forces
$|Z|>n_q-m/2\ge n_q/2$. Chebyshev and a union bound give
\begin{equation}
 \Pr(\mathcal B_q^c)\le4T/n_q\le4T/m.
 \label{frames:stopped-bad}
\end{equation}

We finish by coupling the endpoint lists, conditional on $\mathbf c$.
Write $\pi_q$ for the uniform injection law on $q$ positions, and put
\[
 \delta_{\rm aux}(\mathbf c)
 =\left\|\mathcal L((Y_T(1),\ldots,Y_T(q))\mid\mathbf c)-\pi_q
       \right\|_{\rm TV}.
\]
At stage $j$, prescribe the real endpoint kernel to be $p_{j,T}$,
evaluated on the preceding real paths. This prescription is used
for every joint coupling history; it is not replaced by a law conditioned
on earlier coupling success. Prescribe the comparison kernel to be
uniform on the positions not yet used by the comparison list.

On previous success the two lists have the same preceding endpoints,
so their available endpoint sets agree. If $\mathcal B_j$ holds,
maximally couple the prescribed kernels, declaring failure if their
endpoints differ. If balance fails, declare failure
and use any coupling of them. Once a failure has occurred, continue
using any coupling of the prescribed kernels, with the comparison
list's own remaining set. After choosing the real endpoint, sample its
full real path from the conditional path law given that endpoint,
$\mathbf c$, and the preceding real paths. These prescriptions preserve the
real path law and the uniform comparison-list law at every stage.
The probability of a balance failure before any mismatch is at most
$\Pr(\mathcal B_q^c\mid\mathbf c)$ by nesting. The mismatch probability on
a balanced, previously successful stage is bounded by
$\tfrac12\sqrt{n_jE_{j,T}}$. Averaging over $\mathbf c$ and dropping
the preceding-success indicator gives the total bound
\begin{align*}
 \mathbb E_{\mathbf c}\delta_{\rm aux}(\mathbf c)
 &\le \Pr(\mathcal B_q^c)
 +\frac12\sum_{j=1}^q
       \mathbb E\!\left[\mathbf1_{\mathcal B_j}
                         \sqrt{n_jE_{j,T}}\right]\\
 &\le\frac{4T}{m}
     +\frac12n^{3/2}(1-1/(4b))^{(T-2d+1)/2},
\end{align*}
by \eqref{frames:stopped-energy}, \eqref{frames:stopped-bad}, and
Cauchy--Schwarz. The fixed-list frame comparison identifies the left
side with the physical marginal distance, proving
\eqref{frames:stopped-delta}.
Finally the exponential factor is at most $e^{-5d/2}$, so
$n^{3/2}e^{-5d/2}\to0$, and $T/m\to0$.
\end{proof}

For sixteen equal blocks, these ordinary $o(1)$ complementary marginal
errors also satisfy the common-trimming hypotheses of
Proposition~\ref{shear:prop:fourier-inputs}. In the following application
we use the spectral-pruning theorem of the companion
(Theorem~\Lref{l:spectral-pruning}) for its operator bound with constant
one. If all
sixteen fixed complementary marginals of a probability $\mu_n$ have
TV error $o(1)$, then there is a subprobability $\nu_n\le\mu_n$ of
mass $1-o(1)$ satisfying
$\|\widehat\nu_n(\rho)\|_{\rm op}\le D^{-1/8}$ for every $D>1$.
The theorem takes one union of coefficient tests for each complementary
marginal. Thus its proof does not multiply the marginal error by the
number of tests. Proposition~\ref{frames:stopped-marginals} supplies
exactly these fixed-marginal hypotheses.

\begin{proposition}[The stopped-energy and pruning application]
\label{frames:pruning-application}
The stopped partial-observation estimate and spectral pruning give
full-permutation distance $o(1)$ by time $5152d$.
\end{proposition}

\begin{proof}
Use any fixed partition into sixteen equal blocks and apply
Proposition~\ref{frames:stopped-marginals} to its complements.
Theorem~\Lref{l:spectral-pruning} gives exponent $a=1/8$ with
constant one. Sixteen convolutions have exponent
$2-2(16)/8=-2$. The original block has zero sign mean, and
Proposition~\Lref{l:amplification} applies to its pruned subprobability.
The total number of physical steps is $16(322d)=5152d$.
\end{proof}

\subsection{One history event that works for every tilted law}
\label{frames:entropy-section}

{\interlinepenalty=10000
The preceding routes control partial marginals for the original law.
Here we seek an entropy comparison that remains valid for every law
concentrated on one common set of histories. The comparison will apply
after conditioning on an event chosen from an arbitrary representation
coefficient. Its probabilistic input is a uniform control of the
conditional likelihoods of realized card endpoints.\par}

Assume $d\ge7$, and fix $L=128$, $q=n-n/L$, and
$T=(1+100L)d=12801d$. Let $\mathbb P$ be the law of cyclic switch
histories, with cumulative label-to-position maps $g_t$ and cyclic
directions $w_t$. Before layer $t$, let $M_t$ be the matching of
trajectory labels: $u,v$ form a pair when
$g_{t-1}(v)=g_{t-1}(u)+w_t$.

For a nonempty subset $F$ of labels, of size $m$, let $B_t^F$ be the
matrix on $F$ which averages the two coordinates of every pair of $M_t$
contained in $F$ and fixes the other coordinates. Put
\[
 Q_t^F=B_1^F\cdots B_t^F,\qquad Q_0^F=I_F.
\]
These matrices depend on the realized matching history. The next lemma
identifies their rows with the conditional hidden-card flow, expressed
using that history to label the available positions.

\begin{lemma}[Conditional masses in trajectory-label coordinates]
\label{frames:reference-rows}
Fix a realized cyclic history $g=(g_0,\ldots,g_T)$ and a nonempty set
$F$ of labels. For $u\in F$, let $p_t^{u,F}(x)$ be the conditional
probability under $\mathbb P$ that card $u$ is at $x$ at time $t$, given
the complete trajectories of the labels in $V\setminus F$, evaluated
on the trajectories supplied by $g$. Define a row indexed by $F$ by
\[
 r_t^{u,F}(v)=p_t^{u,F}(g_t(v)),\qquad v\in F.
\]
Then, for this realized history,
\begin{equation}
 r_0^{u,F}=e_u^{\mathsf T},\qquad
 r_t^{u,F}=r_{t-1}^{u,F}B_t^F,
 \qquad r_t^{u,F}=e_u^{\mathsf T}Q_t^F.
 \label{frames:reference-row-flow}
\end{equation}
Here $e_u$ is the coordinate vector in $\mathbb R^F$.
In particular the conditional likelihood of the realized endpoint of
card $u$ is
\begin{equation}
 p_T^{u,F}(g_T(u))=Q_T^F(u,u).
 \label{frames:diagonal-shadow}
\end{equation}
The conditioning in this statement remains on the paths of
$V\setminus F$; the realized history enters $r_t^{u,F}$ through the
indexing bijection $v\mapsto g_t(v)$.
\end{lemma}

\begin{proof}
The conditioned paths are feasible, because they come from $g$.
Their complement at time $t$ is exactly the set of available positions
$g_t(F)$. Lemma~\ref{frames:holes} gives the conditional mass flow on
these positions. We express its three edge cases in the reference
indices.

If a matching edge has one position in $g_{t-1}(F)$, its coin is fixed
by the path of the observed card on that edge. It is therefore the coin
used by the reference history, and it carries the available position
$g_{t-1}(v)$ to $g_t(v)$. Its mass keeps index $v$. If both endpoints
belong to $g_{t-1}(F)$, say with reference indices $v,v'$, the two
conditional output masses are the average of the preceding two masses.
They are equal, so the reference history's choice to swap or fix that
edge does not change the two indexed output values. An edge with no
available position has no coordinate in the row. These cases give
$r_t^{u,F}=r_{t-1}^{u,F}B_t^F$. Initially card $u$ is at $u$, so the
initial row is $e_u^{\mathsf T}$. Iteration proves the row formula, and
the index of the realized endpoint $g_T(u)$ is $u$.
\end{proof}

Let $J_F$ be the matrix all of whose entries equal $1/m$, and set
\[
 Z_t^F=\|Q_t^F-J_F\|_{\rm HS}^2.
\]
By Lemma~\ref{frames:reference-rows}, this is the sum over $u\in F$ of
the squared conditional deviations of card $u$ from uniform on the
available positions, with each deviation computed under the same
observed complement. No independence among those cards is required.
Each $B_t^F$ is a doubly stochastic contraction, hence so is $Q_t^F$.
Since $J_FQ_t^F=J_F$, we have
$Z_t^F\le\|I_F-J_F\|_{\rm HS}^2=m-1\le n$.

For the endpoint likelihood in the lemma, its diagonal entry and the
definition of $Z_T^F$ give
\begin{equation}
 \log\bigl(m p_T^{u,F}(g_T(u))\bigr)
 \le m\sqrt{Z_T^F}.
 \label{frames:likelihood-excess}
\end{equation}
Indeed $Q_T^F(u,u)\le1/m+\sqrt{Z_T^F}$ and
$\log(1+x)\le x$. The likelihood is positive for the feasible realized
history. We will control the average of this logarithmic excess on a
single event, simultaneously for every reservoir size needed by a list.

\begin{lemma}[One typical event for all reservoir sizes]
\label{frames:typical-history}
Define the event on cyclic histories by
\begin{equation}\label{frames:typical-average}
 \mathcal G_{\mathrm{ent}}
 =\bigcap_{m=n/L}^{n}
   \left\{g:\mathbb E_{F:\,|F|=m}Z_T^F(g)\le n^{-8}\right\},
\end{equation}
where each average is uniform over subsets of labels. Then
$\mathbb P(\mathcal G_{\mathrm{ent}})=1-o(1)$.
\end{lemma}

\begin{proof}
We estimate the distribution of these matching functionals using the
directly generated variable-direction process. Start its directions
with the standard basis, and choose each later direction uniformly
outside the span of its preceding $d-1$ directions. At each layer use
fresh independent fair switch coins. Denote its cumulative
permutation by $\widetilde g_t$. Lemma~\ref{frames:matching-invariance} says that,
for each complete direction schedule, its matching sequence on trajectory
labels has the same law as the cyclic matching sequence. Thus every
functional $Z_T^F$ of those matchings has the same law in the two models.
For the following calculation, $B_t^F$, $Q_t^F$, and $Z_t^F$ denote
these same matrix functions evaluated on the auxiliary matching sequence.

Take $F$ uniform of a fixed size $m\ge n/L$, independently of this
auxiliary process. Condition on $F$, the directions through time $t-1$,
and the complete switch history through that time. For a row of
$Q_{t-1}^F-J_F$, write its entries as $z_u$, $u\in F$. This row is
known under the conditioning and has sum zero. Transfer it to the current
positions by
\[
 \widetilde z(\widetilde g_{t-1}(u))=z_u\quad(u\in F).
\]
The label pairs averaged by $B_t^F$ correspond exactly to the position
pairs contained in $\widetilde g_{t-1}(F)$ in the new matching. Thus the squared
row loss is the position-space loss in Lemma~\ref{frames:cross-energy}.

For $t>d$, let $J_t$ be the span of the preceding $d-1$ directions,
and let $h,h'$ be the sizes of
$\widetilde g_{t-1}(F)\cap J_t$ and $\widetilde g_{t-1}(F)\setminus J_t$. The next direction
is fresh under the conditioning above. Applying the lemma to each row
and summing gives
\[
 \mathbb E[Z_{t-1}^F-Z_t^F\mid
       F,\text{ complete generated past through }t-1]
 \ge\frac{\min(h,h')}{n}Z_{t-1}^F.
\]

Now fix the generated past but leave $F$ unobserved. The map $\widetilde g_{t-1}$
and the half-space $J_t$ are fixed, and $\widetilde g_{t-1}(F)$ is a uniform
$m$-subset. The elementary half-count estimate gives
\[
 \Pr\{\min(h,h')<m/4\mid\text{ generated past}\}
 \le 2\theta^m,
 \qquad \theta=(3/2)2^{-3/4}<1.
\]
For example, expansion of the product for a uniform sample without
replacement gives $\mathbb E2^X\le(3/2)^m$ for either half-count $X$;
Markov's inequality applied to the larger half yields this bound.
Set $c=-\log\theta>0$. With $Z_{t-1}^F\le n$ and $m\ge n/L$, the unconditional recurrence is
\[
 \mathbb E Z_t^F\le(1-1/(4L))\mathbb E Z_{t-1}^F
                         +O(ne^{-cn/L})\qquad(t>d).
\]
There are $T-d=100Ld$ useful layers. Starting with $Z_d^F\le n$,
the initial contribution is at most $ne^{-25d}$ and the accumulated
error is exponentially small in $n$. Hence
$\mathbb E Z_T^F\le n^{-20}$ for large $d$, uniformly in $m$.

The equality of matching laws proved at the start transfers this
expectation to the cyclic histories on which $\mathcal G_{\mathrm{ent}}$ was defined.
For each $m$, Markov's inequality applied to the subset average at
threshold $n^{-8}$ gives failure probability at most $n^{-12}$.
There are at most $n$ sizes, so
$\mathbb P(\mathcal G_{\mathrm{ent}}^c)\le n^{-11}=o(1)$.
\end{proof}

For probability laws on a finite set, write
$D(\nu\|\zeta)=\sum_\omega\nu(\omega)\log(\nu(\omega)/\zeta(\omega))$
for relative entropy, with $0\log0=0$ and value $+\infty$ when absolute
continuity fails. Write $\operatorname{Ent}$ for Shannon entropy, using
natural logarithms in both cases.
For a set $X$ of input labels, let $\nu_X$ be the endpoint restriction
law under a law $\nu$ on histories, and let $U_X$ be the uniform injection
law. The typical event now controls the likelihood term in the entropy
chain rule.

\begin{lemma}[Entropy comparison for arbitrary tilted histories]
\label{frames:shadow-entropy}
For every probability law $\nu$ on histories supported on $\mathcal G_{\mathrm{ent}}$,
\begin{equation}\label{frames:entropy-comparison}
 D(\nu\|\mathbb P)\ge
       \mathbb E_{X:\,|X|=q}D(\nu_X\|U_X)-n^{-2}.
\end{equation}
\end{lemma}

\begin{proof}
Fix an ordered list $i_1,\ldots,i_q$. Let $\tau_j$ and $y_j$ be the
complete trajectory and endpoint of $i_j$. For the current history let
$p_j$ be the conditional probability under $\mathbb P$ of its realized
endpoint $y_j$, given $\tau_1,\ldots,\tau_{j-1}$. Relative-entropy
chain rule and conditional data processing give
\begin{align*}
 D(\nu\|\mathbb P)
 &\ge\sum_{j=1}^q\bigl[-\operatorname{Ent}_\nu
     (y_j\mid\tau_1,\ldots,\tau_{j-1})-\mathbb E_\nu\log p_j\bigr]\\
 &\ge-\operatorname{Ent}_\nu(y_1,\ldots,y_q)
                              -\sum_{j=1}^q\mathbb E_\nu\log p_j.
\end{align*}
The second inequality uses that the preceding endpoints are functions
of the preceding trajectories. Put
$F_j=V\setminus\{i_1,\ldots,i_{j-1}\}$ and $m_j=|F_j|=n-j+1$.
The uniform injection assigns each endpoint list mass
$\prod_jm_j^{-1}$, so the last inequality is
\[
 D(\nu\|\mathbb P)\ge D(\nu_X\|U_X)
                     -\sum_{j=1}^q\mathbb E_\nu\log(m_jp_j),
 \qquad X=\{i_1,\ldots,i_q\}.
\]

Lemma~\ref{frames:reference-rows} applies with observed complement
$V\setminus F_j$ and gives $p_j=Q_T^{F_j}(i_j,i_j)$ on each realized
history. For a uniform ordered input list, $F_j$ is a uniform subset
of size $m_j$. The bound \eqref{frames:likelihood-excess} is independent
of which $i_j\in F_j$ is chosen. On every history in $\mathcal G_{\mathrm{ent}}$,
Jensen's inequality therefore gives
\[
 \mathbb E_{\text{ordered list}}\log(m_jp_j)
 \le m_j\mathbb E_{F:\,|F|=m_j}\sqrt{Z_T^F}
 \le n\sqrt{n^{-8}}=n^{-3}.
\]
All these sizes are at least $n/L$. Sum over $q\le n$ indices and
average the preceding entropy inequality over lists and over $\nu$.
The induced set $X$ is uniform of size $q$, and ordering its coordinates
does not change its relative entropy from a uniform injection. This
proves the claimed error $n^{-2}$.
\end{proof}

\begin{proposition}[Representation-coefficient tails from entropy]
\label{frames:entropy-tail}
There is an absolute $C$ such that for every irreducible unitary
representation $\rho$ of $G$, of degree $D$, and unit vectors $u,z$,
\begin{equation}
 \mathbb P\{\mathcal G_{\mathrm{ent}},
       |\langle z,\rho(g_T)u\rangle|\ge D^{-1/8}\}
 \le C D^{-1/(2L)}.
 \label{frames:coefficient-tail}
\end{equation}
Put $a=\mathbb P(\mathcal G_{\mathrm{ent}})$. When $a>0$, the conditional endpoint law
$\eta=\mathcal L(g_T\mid\mathcal G_{\mathrm{ent}})$ is defined and satisfies
\[
 \|\widehat\eta(\rho)\|_{\rm op}
 \le D^{-1/8}+\frac{C}{a}D^{-1/(2L)}.
\]
For all sufficiently large $d$, Lemma~\ref{frames:typical-history}
gives $a\ge1/2$. In that range the last bound is at most
$C'D^{-1/(2L)}$ for the absolute constant $C'=1+2C$.
\end{proposition}

\begin{proof}
If $a=\mathbb P(\mathcal G_{\mathrm{ent}})=0$, the probability on the left side of
\eqref{frames:coefficient-tail} is zero. Otherwise apply the all-tilt
coefficient theorem, \cite[Theorem~\Lref{l:tilted-entropy}]{partial-fourier}.
Its reference probability is the law $\mathbb P$ of cyclic switch
histories, its endpoint map is $g_T$, and its common event is
$\mathcal G_{\mathrm{ent}}$ from Lemma~\ref{frames:typical-history}.
Lemma~\ref{frames:shadow-entropy} verifies the theorem's hypothesis
for every probability $\nu$ supported on that same event, with
$q=n-n/128$ and additive error $n^{-2}$. The theorem therefore gives
the tail exponent $1/256=1/(2L)$ at threshold $D^{-1/8}$.
For $a>0$, splitting each absolute coefficient at $D^{-1/8}$ under
$\mathbb P(\,\cdot\mid\mathcal G_{\mathrm{ent}})$ gives the stated operator bound
with factor $C/a$. By Lemma~\ref{frames:typical-history},
$a=1-o(1)$, so $a\ge1/2$ eventually. Since $L=128$ and $D\ge1$,
$D^{-1/8}\le D^{-1/(2L)}$, which gives $C'=1+2C$.
\end{proof}

\begin{proposition}[The typical-history entropy application]
\label{frames:entropy-application}
With $L=128$, $T=(1+100L)d$, and $b=32L$, the shuffle has
full-permutation distance $o(1)$ at time $bT$.
\end{proposition}

\begin{proof}
For all sufficiently large $d$, put $a=\mathbb P(\mathcal G_{\mathrm{ent}})>0$
and $\eta=\mathcal L(g_T\mid\mathcal G_{\mathrm{ent}})$. The original block law
$\mu=\mathcal L(g_T)$ has zero sign mean, so
\[
 |\widehat\eta(\operatorname{sgn})|
 =\frac{|\mathbb E[\operatorname{sgn}(g_T)\mathbf1_{\mathcal G_{\mathrm{ent}}^c}]|}{a}
 \le\frac{1-a}{a}=o(1).
\]
For the nonlinear irreducibles, Proposition~\ref{frames:entropy-tail}
gives the operator bound with the absolute constant $C'$. The trivial
coefficient of $\eta^{*b}$ is one. Plancherel therefore gives
\[
 |G|\sum_g|\eta^{*b}(g)-U_G(g)|^2
 \le |\widehat\eta(\operatorname{sgn})|^{2b}
       +(C')^{2b}\sum_{\rho\notin\{\mathbf1,\operatorname{sgn}\}}
                   D_\rho^{\,2-b/L}=o(1).
\]
Indeed $2-b/L=-30$, and the nonlinear sum vanishes by
\cite[Lemma~\Lref{l:degree-reciprocal-twenty}]{partial-fourier}.
Cauchy--Schwarz gives total-variation convergence.
The $b$ independent histories all belong to $\mathcal G_{\mathrm{ent}}$ with
probability $a^b$. Removing their conditioning costs at most
$1-a^b\le b(1-a)=o(1)$, because $b$ is fixed.
The physical block length is a multiple of $d$, and its $b$-fold
product is the actual shuffle at time $32L(1+100L)d$.
\end{proof}

\section{An averaged application of deferred columns}
\label{frames:aux:section}

Lemma~\ref{shear:lem:keys} constructs the frame from independent shears,
reads each coefficient when its target column is next selected, and
permits the remaining shear data to be disclosed at the terminal time.
We now give a separate quantitative application of that shared law.
It samples the input ordering and uses simultaneous balance across
all affine half-spaces, then uses the companion's inverse-twentieth
degree estimate and a one-vector orbit bound.

The inverse frame starts at the identity, so its endpoint comparison
preserves the distance for every fixed input list. The averaging in this
application enters only through its balance estimate. Its full-deck
consequence is the following bound.

\begin{theorem}\label{frames:aux:theorem}
For the Thorp shuffle on $n=2^d$ labeled cards,
\[
 \max_{\sigma\in S_n}
 \bigl\|P_d^{512040d}(\sigma,\cdot)-U_n\bigr\|_{\mathrm{TV}}
 =o(1)\qquad(d\longrightarrow\infty),
\]
where $U_n$ is the uniform probability measure on $S_n$ and one time step is
one physical shuffle. Consequently, at total-variation threshold $1/4$,
\[
 d\le t_{\mathrm{mix}}(d)\le512040d
\]
for all sufficiently large integers $d$.
\end{theorem}

Throughout this section put
\begin{equation}\label{frames:aux:parameters}
 r=128,\qquad k=(1-1/r)n,\qquad T=(100r+1)d.
\end{equation}
We take $d\ge7$, so $r$ divides $n$, and let $\mu$ be the coordinate
shuffle law after $T$ layers. We first bound the average marginal distance
for ordered lists of $k$ cards. Forty independent blocks of this length
will then give the theorem.

\subsection{The inverse frame and its conditional law}

Let $\iota_1,\iota_2,\ldots$ be a deterministic periodic sequence of
coordinate indices, with period $d$ containing each of $1,\ldots,d$
once. For the present application use the standard cyclic order. For
each layer $i$, choose an independent uniform linear functional
$\ell_i:V\to\mathbb F_2$ satisfying $\ell_i(e_{\iota_i})=0$, and put
\begin{equation}\label{frames:aux:transvections}
 S_i x=x+\ell_i(x)e_{\iota_i},\qquad
 B_0=I,\quad B_i=S_iB_{i-1},\qquad
 w_i=B_{i-1}^{-1}e_{\iota_i}.
\end{equation}
Each $S_i$ is an involution. Write
$\Lambda=(\ell_1,\ldots,\ell_T)$ and $W=(w_1,\ldots,w_T)$.
Conditional on $\Lambda$, let $M_i$ be independent uniform matching
switches in directions $w_i$, and define
\[
 Y_0=I,\qquad Y_i=M_iY_{i-1},\qquad X_i=B_iY_i.
\]
The process $Y$ uses virtual positions and $X$ uses the original
coordinate positions.

\begin{lemma}[Inverse-frame form of the shear law]
\label{frames:aux:directions}
Conditional on $\Lambda$, the increments of $X$ are independent
uniform switches in the coordinate directions $e_{\iota_i}$. For every
$t\le T$, the conditional law of $(Y_i)_{i=0}^t$ given $\Lambda$ is
the product-switch law determined by $w_1,\ldots,w_t$.

Every $d$ consecutive directions $w_i$ form a basis. At layer $i$,
learning $w_i$ does not change the conditional law of the preceding
virtual positions given $w_1,\ldots,w_{i-1}$ and any specified virtual
path prefixes through layer $i-1$. For $i>d$, that direction is uniform
outside
\[
 \operatorname{span}(w_{i-d+1},\ldots,w_{i-1})
\]
conditional on that same direction and path-prefix field. Finally,
conditional on the complete
direction sequence $W$, disclosing the remaining information in
$\Lambda$ does not change the conditional law of the virtual
trajectories, including after a fixed collection of those trajectories
has been observed.
\end{lemma}

\begin{proof}
Set $A_i=B_i^{-1}$. Since $S_i$ is an involution,
$A_i=A_{i-1}S_i$ and $w_i=A_{i-1}e_{\iota_i}$, which are the forward
frame and selected directions of Lemma~\ref{shear:lem:keys}.
Remark~\ref{shear:rem:factorization} gives the inverse coupled identity
$X_i=B_iY_i$ and the conditional coordinate-switch law. The prefix
product-switch law and consecutive-basis property are conclusions of
that lemma for the same periodic axis order. The prefix law makes the
old virtual path and the next direction conditionally independent given
the direction prefix; Corollary~\ref{frames:conditional-interface}
records the resulting hidden-position factorization after specified
virtual path prefixes through layer $i-1$ are observed. The
uniform-complement law for $i>d$ is
the lemma's fresh-direction conclusion.

In particular, the conditional law of the full virtual trajectory
given $\Lambda$ depends on $\Lambda$ only through $W$. Thus the
trajectory and $\Lambda$ are conditionally independent given $W$.
Conditioning also on a fixed observed part of the trajectory preserves
this factorization for the remaining trajectories.
\end{proof}

For the standard cyclic order used in this application, fix an input
list and $\Lambda$. The known output
bijection $B_T$ preserves uniform injections and total variation.
The conditional coordinate law in the lemma therefore identifies that
list's virtual conditional distance with its distance under $\mu$.
The random ordering introduced next is used to prove balance; it is not
required by this frame comparison.

\subsection{A conditional energy estimate for almost all cards}

Choose a uniformly random ordering
$\mathcal O=(a_1,\ldots,a_n)$ of the cards, independently of
$(\Lambda,(Y_i)_{i=0}^T)$. For $0\le j<k$, use the chronological field
\[
 \mathcal F_i^j
 =\sigma\bigl(\mathcal O,w_1,\ldots,w_i,\,
       Y_s(a_h):0\le s\le i,\ 1\le h\le j\bigr).
\]
Let $D_i=V\setminus\{Y_i(a_h):1\le h\le j\}$ be the available
positions, of size $m=n-j$. Define
\[
 q_i(x)=\Pr\{Y_i(a_{j+1})=x\mid\mathcal F_i^j\},
\]
and put
\[
 z_i(x)=q_i(x)-\frac1m\quad(x\in D_i),\qquad
 E_i=\sum_{x\in D_i}z_i(x)^2.
\]
The conditional mass is zero off $D_i$, and $z_i$ sums to zero on
$D_i$. We use counting measure in the squared norm, so
$E_0=1-1/m\le1$.

\begin{lemma}\label{frames:aux:energy}
Uniformly for $0\le j<k$, with expectation over the ordering and shuffle
construction,
\begin{equation}\label{frames:aux:energy-bound}
 \mathbb E E_T\le e^{-25d}+T\delta_n,
 \qquad
 \delta_n=2n(n+1)\exp\left(-\frac{n}{16r}\right).
\end{equation}
\end{lemma}

\begin{proof}
Lemma~\ref{frames:aux:directions} gives the conditional product-switch
law and the fresh direction in the field $\mathcal F_{i-1}^j$.
Apply Lemma~\ref{shear:lem:energy} with initial datum $\mathcal O$,
$t=i-1$, available set $D_{i-1}$, and hidden mass $q_{i-1}$.
Its averaging and transport rules preserve the uniform comparison law
on the holes and give
\begin{equation}\label{frames:aux:energy-drop}
 E_{i-1}-E_i
 =\frac12\sum_{\substack{\{x,y\}\text{ an edge of layer }i\\
                         x,y\in D_{i-1}}}
       \bigl(z_{i-1}(x)-z_{i-1}(y)\bigr)^2.
\end{equation}
In particular $0\le E_i\le E_{i-1}\le1$ for every realization.

Suppose now that $i>d$. Before its direction is drawn, put
$J_i=\operatorname{span}(w_{i-d+1},\ldots,w_{i-1})$ and split
$D_{i-1}=D^0\cup D^1$ across its two cosets, with sizes $h_0,h_1$.
Lemma~\ref{frames:aux:directions} makes the next direction uniform
outside $J_i$, so each cross pair has conditional matching probability
$2/n$. The cross-pair calculation in
Lemma~\ref{shear:lem:energy}, with centered vector $z_{i-1}$, gives
\begin{align}
 \mathbb E(E_{i-1}-E_i\mid\mathcal F_{i-1}^j)
 &=\frac1n\sum_{x\in D^0,\ y\in D^1}
             \bigl(z_{i-1}(x)-z_{i-1}(y)\bigr)^2\notag\\
 &\ge\frac{\min(h_0,h_1)}n E_{i-1}.
 \label{frames:aux:cross-loss}
\end{align}
The inequality is \eqref{shear:eq:bipartite-form}; it also covers an
empty part.

We use a simultaneous balance estimate for every affine half-space.
At a fixed time, condition on $\Lambda$ and the complete virtual switch
history but not on $\mathcal O$. The permutation $Y_{i-1}$ is then fixed,
and the independent ordering makes $D_{i-1}$ a uniform $m$-subset of
$V$. The simultaneous estimate below is valid even for a half-space
chosen after that subset is seen.

For $m<n$, a uniform $m$-subset has the law of independent
Bernoulli choices of parameter $m/n$, conditional on making exactly $m$
choices.  This conditioning event has probability at least $1/(n+1)$:
$m$ is a mode of the binomial distribution and that distribution has
$n+1$ values. In a specified affine half-space let $X$ be the
unconditioned number of choices. It has mean $m/2$, and
\begin{align*}
 \Pr\{X<m/4\}
 &\le 2^{m/4}\mathbb E2^{-X}
 =2^{m/4}\left(1-\frac{m}{2n}\right)^{n/2}\\
 &\le e^{-(1-\log2)m/4}\le e^{-m/16},
\end{align*}
where the last inequality uses $\log2<3/4$. There are
fewer than $2n$ affine half-spaces arising as a linear hyperplane or its
other coset.  Since $m\ge n/r$, a union bound after conditioning gives
\begin{equation}\label{frames:aux:balance}
 \mathbb P\bigl(\text{some such half-space contains fewer than }m/4
                  \text{ holes}\bigr)\le\delta_n.
\end{equation}
For $m=n$ the balance assertion is deterministic.

On the complement of the event in \eqref{frames:aux:balance}, the
conditional loss in \eqref{frames:aux:cross-loss} is at least
$mE_{i-1}/(4n)\ge E_{i-1}/(4r)$.  On the exceptional event we retain
the deterministic energy nonincrease and the bound $E_{i-1}\le1$.
Taking expectations consequently gives
\[
 \mathbb E E_i\le
 \left(1-\frac1{4r}\right)\mathbb E E_{i-1}+\delta_n
 \qquad(i>d).
\]
Only the unconditional estimate \eqref{frames:aux:balance} has been
used.  Iterating over $T-d=100rd$ layers, starting from $E_d\le1$,
yields
\[
 \mathbb E E_T\le
 \left(1-\frac1{4r}\right)^{100rd}+T\delta_n
 \le e^{-25d}+T\delta_n.
\]
\end{proof}

The terminal disclosure and the common tuple comparison now convert
this energy estimate to the required average of marginal distances.

\begin{proposition}\label{frames:aux:marginal}
Let $C=(a_1,\ldots,a_k)$ be a uniformly chosen ordered list of distinct
initial cards, independent of $X_T\sim\mu$, and let $\pi_k$ be uniform
on ordered $k$-tuples of distinct positions. Then
\begin{align}
 \mathbb E_C\bigl\|\mathcal L(X_T(C)\mid C)-\pi_k\bigr\|_{\rm TV}
 &\le\varepsilon_n,\notag\\
 \varepsilon_n
 &=\frac{k}{2}\sqrt{n\bigl(e^{-25d}+T\delta_n\bigr)}=o(1).
 \label{frames:aux:marginal-error}
\end{align}
\end{proposition}

\begin{proof}
For each $j<k$, the terminal factorization in
Lemma~\ref{frames:aux:directions} and
\eqref{frames:terminal-disclosure} give
\[
 \Pr\{Y_T(a_{j+1})=x\mid
       \Lambda,\mathcal O,\,
       Y_s(a_h):0\le s\le T,\ h\le j\}=q_T(x).
\]
Thus the terminal-data hypothesis of Lemma~\ref{shear:lem:tuple} holds
with initial datum $\mathcal O$ and terminal datum $\Lambda$. Write
\[
 \delta_{\rm virt}(\Lambda,\mathcal O)
 =\left\|\mathcal L\bigl((Y_T(a_1),\ldots,Y_T(a_k))
              \mid\Lambda,\mathcal O\bigr)-\pi_k\right\|_{\rm TV}
\]
for the conditional virtual tuple distance. The lemma's coarsening to
preceding endpoints and sequential comparison, together with
Lemma~\ref{frames:aux:energy}, give
\[
 \mathbb E_{\Lambda,\mathcal O}\delta_{\rm virt}(\Lambda,\mathcal O)
 \le\frac{k}{2}\sqrt{n\bigl(e^{-25d}+T\delta_n\bigr)}.
\]
For each fixed $\Lambda$ and $\mathcal O$, applying $B_T$ to every
output preserves $\pi_k$ and total variation. The conditional law of
$X_T$ given $\Lambda$ is $\mu$, and the independent ordering contributes
no further information about that law. Hence
$\delta_{\rm virt}(\Lambda,\mathcal O)$ equals the distance under
$\mu$ for the specified list $C$.
Averaging proves \eqref{frames:aux:marginal-error}. Finally, $n=2^d$ gives
$n^{3/2}e^{-25d/2}\to0$, while $\delta_n$ is exponentially small in
$n$ up to a polynomial factor.  Hence $\varepsilon_n=o(1)$.
\end{proof}

For the representation step, write
$K_0=\sup_{s\ge1}\sum_{\rho\in\operatorname{Irr}(S_s)}D_\rho^{-20}$.
The proof of Lemma~\Lref{l:degree-reciprocal-twenty} gives $K_0<\infty$,
and the lemma gives vanishing of the same sum after omitting the trivial
and sign types as $s\to\infty$. In particular, for every $M\ge1$,
\[
 \#\{\rho\in\operatorname{Irr}(S_s):D_\rho\le M\}\le K_0M^{20}.
\]
This is the hook-product version of the inverse-degree estimate in the
companion; its proof bounds the first-row hooks by opposite-order
rearrangement. We use this power-twenty estimate in the orbit bound below.

\subsection{Coset truncation and the orbit of one vector}

Choose a uniform partition of the card labels into $r$ blocks
$C_1,\ldots,C_r$ of equal size $n/r$.  For each block its complement
is a uniform $k$-element subset.  Ordering that complement in any
fixed way identifies its labeled endpoint images with an ordered
tuple and does not affect total-variation distance.  By
Proposition~\ref{frames:aux:marginal}, the expected sum of the $r$
complement marginal distances is at most $r\varepsilon_n$.  Fix a
partition realizing a sum at most this bound, and let $H_i\cong
S_{n/r}$ be the subgroup of $S_n$ permuting $C_i$ and fixing its
complement pointwise.

Two position permutations agree on the complement of $C_i$ exactly
when they lie in the same left coset $gH_i$.  Thus the corresponding
coset marginal of $\mu$ is within its complement marginal distance
of uniform on $G/H_i$, where $G=S_n$.  Delete from $\mu$ every
element lying in a coset, for any $i$, whose $\mu$-mass exceeds
$2/[G:H_i]$, and call the resulting subprobability measure $\nu$.
Writing $p=\nu(G)$, we obtain
\begin{equation}\label{frames:aux:cosets}
 1-p\le2r\varepsilon_n=o(1),\qquad
 \nu(gH_i)\le\frac2{[G:H_i]}\quad(g\in G,\ 1\le i\le r).
\end{equation}
Indeed, if a coset has mass $a>2u$, where $u=1/[G:H_i]$, then
$a\le2(a-u)$.  Its removed mass is bounded by twice its positive
excess over uniform; summing positive excesses gives exactly the
marginal total-variation distance.  A union bound handles all $r$
families of deleted cosets.

For a unitary representation $\rho$, define its Fourier transform by
\[
 \widehat\nu(\rho)=\sum_{g\in G}\nu(g)\rho(g).
\]
The following estimate explains why information about complements
of the blocks is enough to control every irreducible representation.

\begin{lemma}\label{frames:aux:orbit-lift}
If $\rho$ is an irreducible unitary representation of $G=S_n$ of
dimension $D$, then the measure in \eqref{frames:aux:cosets} satisfies
\begin{equation}\label{frames:aux:operator-bound}
 \|\widehat\nu(\rho)\|_{\mathrm{op}}
 \le C_rD^{-1/2+11/r},\qquad C_r=\sqrt{2rK_0}.
\end{equation}
\end{lemma}

\begin{proof}
The proof of Lemma~\Lref{l:degree-reciprocal-twenty} bounds the complete
inverse-twentieth sum by $K_0$. Apply Proposition~\Lref{l:orbit-span}
with $p=20$, $K=2$, and $r=128$. That proposition assigns a component
$v_i$ of a given vector to the small $H_i$-types and sets
$W_i=\operatorname{span}\{\rho(h)v_i:h\in H_i\}$. Its orbit bound is
\begin{equation}\label{frames:aux:orbit-dimension}
 \dim W_i
 \le\sum_{\substack{\tau\in\operatorname{Irr}(H_i)\\
                         \dim\tau\le D^{1/r}}}(\dim\tau)^2
 \le K_0D^{22/r}.
\end{equation}
The sum counts each irreducible subgroup type once: the orbit of a
fixed vector in all copies of a type of degree $f$ has dimension at
most $f^2$. Combining this estimate with the coset cap and Schur
averaging in that proposition gives
$\sqrt{2rK_0}D^{-1/2+11/r}$.
\end{proof}

\subsection{Forty blocks and the full-deck conclusion}

The preceding estimate controls all non-linear irreducible types.
We now sum their contributions and treat the sign representation
separately.  Let $b=40$.  Since $0\le\nu\le\mu$, their convolution
powers satisfy $0\le\nu^{*b}\le\mu^{*b}$ and
\begin{equation}\label{frames:aux:lost-mass}
 \|\mu^{*b}-\nu^{*b}\|_1=1-p^b=o(1).
\end{equation}
Here the $\ell^1$ norm uses counting measure on $G$. Apply the centered
identity in Lemma~\ref{lem:subprobability-fourier} to
$\xi=\nu^{*b}$, of mass $p^b$, and then Cauchy--Schwarz. With the
convolution convention of Section~\ref{sec:completion}, this gives
\begin{equation}\label{frames:aux:plancherel}
 \|\nu^{*b}-p^bU_n\|_1^2
 \le\sum_{\substack{\rho\in\operatorname{Irr}(G)\\
                     \rho\ne\mathbf1}}
       D_\rho\|\widehat\nu(\rho)^b\|_{\mathrm{HS}}^2.
\end{equation}
The centering by $p^bU_n$ cancels exactly the trivial coefficient.

For every type other than trivial and sign,
$\|A^b\|_{\mathrm{HS}}\le\sqrt D\|A\|_{\mathrm{op}}^b$,
so Lemma~\ref{frames:aux:orbit-lift} bounds its total contribution
to \eqref{frames:aux:plancherel} by
\[
 C_r^{2b}
 \sum_{\rho\ne\mathbf1,\operatorname{sgn}}
 D_\rho^{\,2-b(1-22/r)}.
\]
For our constants,
\[
 2-40(1-22/128)=-\frac{249}{8}=-31.125<-20.
\]
Lemma~\Lref{l:degree-reciprocal-twenty} thus makes this sum $o(1)$.
No normality of $\widehat\nu(\rho)$ is needed for this norm bound.

Since $d$ divides $T$, the block has no remaining rotation and
$\mu=q_d^{*T}$. For this physical block law,
Lemma~\ref{lem:fair-sign} gives $\widehat\mu(\operatorname{sgn})=0$
and $|\widehat\nu(\operatorname{sgn})|\le1-p$.
The remaining sign contribution in
\eqref{frames:aux:plancherel} is consequently at most
$(1-p)^{2b}=o(1)$.  We conclude from
\eqref{frames:aux:plancherel} that
$\|\nu^{*b}-p^bU_n\|_1=o(1)$.  Combining this with
\eqref{frames:aux:lost-mass} and
$\|p^bU_n-U_n\|_1=1-p^b$ gives
\[
 \|\mu^{*40}-U_n\|_1=o(1).
\]

Each block has length $T=(100r+1)d$, a multiple of $d$.
Independent consecutive blocks therefore have precisely the
convolution law used above, with no unaccounted coordinate rotation.
Since
\[
 40T=40(100\cdot128+1)d=512040d,
\]
this proves the upper assertion of
Theorem~\ref{frames:aux:theorem}.  For any initial labeled deck,
applying a position permutation gives a bijection between $S_n$ and
the possible resulting decks.  It preserves uniform measure and
total-variation distance, which proves the stated worst-start
quantifier.

The matching lower bound is the support calculation in
Remark~\ref{frames:lower-import}. This completes
\mbox{Theorem}~\ref{frames:aux:theorem}.
